\documentclass[11pt]{article}
\usepackage[margin=1.1in]{geometry}
\usepackage{amsmath,amssymb,amsthm}
\usepackage[hidelinks]{hyperref}
\newtheorem{theorem}{Theorem}
\newtheorem{lemma}[theorem]{Lemma}
\newtheorem{corollary}[theorem]{Corollary}
\theoremstyle{remark}
\newtheorem{remark}[theorem]{Remark}
\newcommand{\Col}{\mathrm{Col}}
\newcommand{\Syr}{\mathrm{Syr}}
\newcommand{\dlow}{\underline{d}}
\newcommand{\dup}{\overline{d}}

\title{A zero-or-gap principle for the Collatz exceptional set,\\ and the exponent-one barrier}
\author{Working note (Collatz project), 22 November 2026}
\date{Draft, not refereed. It relies on an unrefereed 2026 preprint.}

\begin{document}
\maketitle

\begin{abstract}
Let $E$ be the set of positive integers whose Collatz orbit never reaches $1$. Assume the 2026
claim that every target $z$ with $3\nmid z$ has a predecessor basin of positive lower density
$d_z$. We observe that $E$ is either empty or has positive lower density, at least
$d_{\Syr(\min E)}$. As a result, Collatz holds as soon as the density $f(y)$ of integers whose
orbit stays $\ge y$ falls below the basin floor $D(1.5y+1)$, for every $y>2^{71}$. Numerically,
and in a branching heuristic, $D(Z)=Z^{-1-o(1)}$. Probabilistic methods give
$f(y)\ll(\log y)^{-c}$. So improving Tao's rate past exponent one, to $f(y)=o(y^{-1-\varepsilon})$,
is equivalent to the conjecture, up to a finite verification. Anything weaker is purely
quantitative.
\end{abstract}

\section{Setting}
$C(n)=n/2$ for even $n$ and $3n+1$ for odd $n$. $\Syr(n)=(3n+1)/2^{\nu_2(3n+1)}$ for odd $n$.
$\Col_{\min}(N)=\min_{k\ge0}C^k(N)$. For $A\subseteq\mathbb N$, $\dlow(A)$ and $\dup(A)$ denote
lower and upper natural density. Define
\[
E=\{N:\ 1\notin\{C^k(N)\}_{k\ge0}\},\qquad
P(z)=\{N:\ z\in\{C^k(N)\}_{k\ge0}\},\qquad
d_z=\dlow(P(z)),
\]
\[
f(y)=\dup\{N:\Col_{\min}(N)\ge y\},\qquad D(Z)=\min\{d_z:\ z\le Z,\ 3\nmid z\}.
\]
For $y\ge2$ we have $\{N:\Col_{\min}(N)\ge y\}\subseteq E$, since an orbit that reaches $1$ has minimum $1$.

\paragraph{Inputs.}
\begin{itemize}
\item[(V)] Every $N<2^{71}$ reaches $1$ (Barina 2025).
\item[(P)] For every $z$ with $3\nmid z$, $d_z>0$ (Shaik, hexagon 2610.00523, Thm.~I.2; also
  Mazur 2026). Both are unrefereed. The bound in (P) depends on $z$; no uniform bound is
  possible, because distinct basins are disjoint.
\end{itemize}

\section{The principle}
\begin{lemma}[zero-or-gap]\label{lem:gap}
Assume (V) and (P). If $E\neq\emptyset$, let $y^*=\min E$. Then $y^*$ is odd, $y^*>2^{71}$,
and $z:=\Syr(y^*)$ satisfies $3\nmid z$ and $z\le(3y^*+1)/2$. Moreover
\[
P(z)\subseteq\{N:\Col_{\min}(N)\ge y^*\}\subseteq E,
\qquad\text{hence}\qquad
f(y^*)\ \ge\ \dlow(E)\ \ge\ d_z\ \ge\ D\big(\tfrac{3y^*+1}{2}\big)>0 .
\]
\end{lemma}
\begin{proof}
If $y^*$ were even, the orbit of $y^*/2$ would pass through $y^*$. Then $y^*/2\in E$, which
contradicts minimality. So $y^*$ is odd, and $y^*\ge2^{71}$ by (V). The value $3y^*+1$ is
$\equiv1\pmod 3$, so $3\nmid z$. Also $z\le(3y^*+1)/2$.

The orbit of $y^*$ passes through $z$, so $z\in E$. Let $N\in P(z)$. Every element $w$ of the
orbit of $N$ either precedes $z$, in which case its orbit contains $z$, or follows $z$, in which
case it lies on the orbit of $z$. In both cases the orbit of $w$ avoids $1$, so $w\in E$ and
$w\ge y^*$. Hence $\Col_{\min}(N)\ge y^*$. The remaining inequalities follow from (P), from
$\dlow\le\dup$, and from monotonicity.
\end{proof}

\begin{corollary}[density race]\label{cor:race}
Assume (V) and (P). The Collatz conjecture holds if and only if
\[
f(y)\ <\ D\big(\tfrac{3y+1}{2}\big)\qquad\text{for every odd } y>2^{71}. \tag{R}
\]
\end{corollary}
\begin{proof}
If Collatz holds, then $E=\emptyset$ and $f\equiv0<D$. Conversely, if $E\neq\emptyset$, then
Lemma~\ref{lem:gap} at $y=y^*$ contradicts (R).
\end{proof}

\begin{remark}[log density]
Tao's theorem bounds the upper \emph{logarithmic} density of $\{\Col_{\min}>N_0\}$ by
$C(\log N_0)^{-c}$. Lower natural density is at most lower logarithmic density, so the lemma
also gives $d_z\le C(\log y^*)^{-c}$. This is consistent, and it is why no contradiction arises
at present.
\end{remark}

\section{Size of the floor, and the barrier}
\paragraph{Numerics.} We counted $\#(P(z)\cap[1,X])$ by depth-first search on the inverse tree,
with $X=10^8$ and $10^9$; the code is \texttt{tools/density/predcount.c}. The quantity $z\,d_z$
is stable in $X$, and over random $z\in[10^4,2\cdot10^4]$ with $3\nmid z$ it has median $\approx7$.
Its left tail satisfies $\Pr(z\,d_z<w)\approx w^{2.9}$ for $w\lesssim1.6$. This is consistent
with $z\,d_z$ being the limit of a branching martingale. Heuristically it gives
$D(Z)\approx Z^{-1-1/2.9}=Z^{-1.35}$.

\paragraph{Proven floor.} Shaik's construction has score $W_z\asymp1/z$. The polynomial mixing
rate then turns this into $d_z\approx\exp(-c z^{0.011})$. If the Syracuse characteristic
function decays like $e^{-In}$, as numerics suggest with $I=I(\log_2 3)\approx0.055$, a modified
bookkeeping gives $d_z\ge c_0 z^{-B}$ with $B\approx3+4.15/I\approx76$. This is our own
derivation and is unrefereed.

\paragraph{The barrier.} Corollary~\ref{cor:race} shows that any proof of $f(y)=o(D(1.5y))$ is
a proof of Collatz. With $D(Z)=Z^{-1-o(1)}$, the threshold for the rate exponent is $1$. If
$f(y)\ll y^{-A}$ with $A<1$, this is consistent with a hypothetical cycle whose basin has density
$\asymp1/y^*$, so such bounds are purely quantitative. No probabilistic method that also applies
to $3x-1$ can reach $A>1$ without new deterministic input.

\paragraph{Kill tests.}
\begin{itemize}
\item $3x-1$: its cycles have minima $1,5,17$, so the analogue of (R) only concerns $y>17$. It
  is consistent.
\item $5x+1$: Tao's method fails there because the drift is positive.
\item The argument uses only (V), (P) and the orbit structure. The arithmetic of $3$ enters
  through (P), via the obstruction for targets with $3\mid z$.
\end{itemize}

\end{document}
